\documentclass[a4paper, 12pt, final]{article}

\usepackage[T2A]{fontenc}
\usepackage[utf8]{inputenc}
\usepackage[english, russian]{babel}

% [showframe] - useful option
\usepackage[margin=2cm]{geometry}

% Установка межстрочного интервала
\usepackage{setspace}
\onehalfspacing

% Установка интервала между абзацами
\setlength{\parskip}{1em}

% Отступ (красная строка) для первого абзаца главы или параграфа
\usepackage{indentfirst}

\usepackage{amsmath}
\usepackage{amsfonts} % шрифты в формулах
\usepackage{amsthm}   % теоремы

% dcases
\usepackage{mathtools}

\newtheorem{definition}{Определение}

% Оглавление и ссылки как гиперссылки
% Это для удобства чтения в электронном виде
\usepackage[unicode, breaklinks]{hyperref}

\usepackage{xcolor}

\begin{document}

\begin{center}
    \scshape
    \hrule
    НЕПРЕРЫВНОСТЬ ФУНКЦИИ
    \vspace{0.7em} \hrule
\end{center}

\begin{flushleft}
    Serge Beloff \& Chat GPT
    \vspace{0.7em} \hrule
\end{flushleft}

\tableofcontents

\section{Аннотация}

Из курса школьной математики известны такие свойства функции $y=f(x)$ как область определения $D$ и область значений $E$.
Например, для функции $f(x)=\frac{1}{x}$ областью определения $D(f)$ будут все вещественные числа кроме нуля: $D(f) = (-\infty; 0) \cup (0; +\infty)$.
Такая же и область значений: $E(f) = (-\infty; 0) \cup (0; +\infty)$.
Для параболы $f(x)=x^2$: $D(f)=\mathbb{R}$, $E(f)=[0; +\infty)$.

Еще одним свойством функции $y=f(x)$ является \textbf{непрерывность}.
Рассмотрим это свойство, докажем непрерывность функции синуса $f(x)=\sin(x)$ и укажем приложения и практический смысл этого свойства.

Интуитивно свойство непрерывности означает плавное (постепенное) изменение значения функции в зависимости от плавного (постепенного) изменения аргумента.
Например, рассмотрим прямолинейное движение поезда, или, вообще материальной точки.
Пусть $x(t)$ --- это уравнение движения.
Мы понимаем, что расстояние $x$ будет изменяться плавно (постепенно).

\section{Определение непрерывности}

\begin{definition}
    Приращением функции $y=f(x)$ в точке $x_0$ называется разность:
    \begin{equation}
        \Delta y = f(x_0 + \Delta x) - f(x_0),
        \label{eq:delta_y}
    \end{equation}
    где $\Delta x$ --- приращение аргумента.
    \label{def:delta_y}
\end{definition}

\begin{definition}
    Функция $y=f(x)$ называется непрерывной в точке $x_0$, если эта функция определена в какой-нибудь окрестности точки $x_0$ и если
    \begin{equation}
        \lim_{\Delta x \to 0} \Delta y = 0,
        \label{eq:cont_def_1}
    \end{equation}
    т.е. если бесконечно малому приращению аргумента $\Delta x$ соответствует бесконечно малое приращение функции $\Delta y$.
    \label{def:cont_def_1}
\end{definition}

\begin{definition}
    Функция $y=f(x)$ называется непрерывной в интервале $[a;b]$, если она непрерывна в каждой его точке.
    \label{def:cont_def_interval}
\end{definition}
Здесь необходимо сделать уточнение для концов интервала.
В точке $x=a$ при вычислении предела мы можем оставаться внутри интервала подходя только справа: $\Delta x \to 0+$.
Аналогично к точке $x=b$ мы можем двигаться только слева: $\Delta x \to 0-$.
Таким образом, функция $y=f(x)$ непрерывна на интервале $[a;b]$, если:
\begin{equation}
    \begin{dcases}
        \lim_{\Delta x \to 0 } &\Delta y = 0, \quad x \in (a;b),\\
        \lim_{\Delta x \to 0+} &\Delta y = 0, \quad x = a,      \\
        \lim_{\Delta x \to 0-} &\Delta y = 0, \quad x = b,
    \end{dcases}
    \label{eq:cont_def_inerval_1}
\end{equation}

Формулу \eqref{eq:cont_def_1} можно записать следующим образом:
\begin{align*}
    \lim_{\Delta x \to 0} \Delta y = 0 \\
    \lim_{\Delta x \to 0} \left[ f(x_0 + \Delta x) - f(x_0) \right] = 0
\end{align*}
Пусть $x=x_0 + \Delta x$.
Тогда при $\Delta x \to 0$, имеем $x \to x_0$.
Значит,
\begin{equation*}
    \lim_{x \to x_0} \left[ f(x) - f(x_0) \right] = 0
\end{equation*}
Выражение $f(x_0)$ не зависит от $x$.
Т.е. является константой.
Поэтому его можно вынести за знак передела:
\begin{equation}
    \lim_{x \to x_0} f(x) = f(x_0)
    \label{eq:cont_def_2}
\end{equation}
В итоге мы получили еще одно определение непрерывности функции:
\begin{definition}
    Функция $y=f(x)$ непрерывна в точке $x_0$, если она определена в какой-нибудь окрестности этой точки и если предел функции при стремлении независимой переменной $x$ к $x_0$ существует и равен значению функции при $x=x_0$.
    \label{def:conf_def_2}
\end{definition}

Из этого следует, что совокупность формул \eqref{eq:cont_def_inerval_1}, описывающая непрерывность функции в интервале, может быть эквивалентно записана следующим образом:
\begin{equation}
    \begin{dcases}
        \lim_{x \to x_0} &f(x) = f(x_0), \quad x_0 \in (a;b),\\
        \lim_{x \to a+ } &f(x) = f(a),                     \\
        \lim_{x \to b- } &f(x) = f(b),
    \end{dcases}
    \label{eq:cont_def_inerval_2}
\end{equation}

\section{Непрерывность синуса}

Докажем непрерывность функции синуса $f(x)=\sin(x)$.
Возможно у Вас, дрогой читатель, есть соблазн воспользоваться определением \ref{def:conf_def_2} и формулой \eqref{eq:cont_def_2}:
\begin{equation*}
    \lim_{x \to x_0} \sin(x) = \sin(x_0)
\end{equation*}
И на том и написать ЧТД.
Но этот путь ложный!
Дело в том, что такая подстановка предполагает, что непрерывность функции $\sin(x)$ уже доказана.
То есть если непрерывность функции доказана, тогда действительно можно смело применять формулу \eqref{eq:cont_def_2}.

Ну что же, пойдем другим путём и сделает настоящее ТРУ доказательство!
Для этого воспользуемся определением \ref{def:cont_def_1} и формулой \eqref{eq:cont_def_1}:
\begin{align*}
    \lim_{\Delta x \to 0} \Delta y =
    \lim_{\Delta x \to 0} \left( \sin(x_0 + \Delta x) - \sin(x_0) \right) = \dots
\end{align*}
Применим формулу синуса суммы:
\begin{equation*}
    \sin (\alpha + \beta) = \sin \alpha \cos \beta + \cos \alpha \sin \beta
\end{equation*}
\begin{align*}
    &= \lim_{\Delta x \to 0} \left( \sin(x_0) \cos(\Delta x) + \cos (x_0) \sin(\Delta x) - \sin(x_0) \right) = \\
    &= \lim_{\Delta x \to 0} \sin(x_0) \cos(\Delta x) +
       \lim_{\Delta x \to 0} \cos(x_0) \sin(\Delta x) -
       \lim_{\Delta x \to 0} \sin(x_0)
\end{align*}
С последним слагаемым всё понятно.
Оно не зависит от $\Delta x$.
То есть:
\begin{equation}
    \lim_{\Delta x \to 0} \Delta y =
    \sin(x_0) \lim_{\Delta x \to 0} \cos(\Delta x) +
    \cos(x_0) \lim_{\Delta x \to 0} \sin(\Delta x) -
    \sin(x_0)
\end{equation}

И вот здесь мы подходим к следующим, так сказать, фундаментальным пределам:
\begin{equation}
    \lim_{x \to 0} \cos(x) = 1
    \label{eq:lim_cos_zero}
\end{equation}
\begin{equation}
    \lim_{x \to 0} \sin(x) = 0
    \label{eq:lim_sin_zero}
\end{equation}
Можно рассмотреть единичную тригонометрическую окружность.
На ней пределы \eqref{eq:lim_cos_zero} и \eqref{eq:lim_sin_zero} интуитивно понятны.

\textcolor{red}{
    TODO: Строгие доказательства пределов \eqref{eq:lim_cos_zero} и \eqref{eq:lim_sin_zero}. Язык $\varepsilon$--$\delta$.
}

\section{Приложения непрерывности}

\phantomsection
\addcontentsline{toc}{section}{Список литературы}

\nocite{Bermant}
\nocite{Fikhtengolts}

\bibliographystyle{apalike}
\bibliography{bib}

\end{document}
